\honest{High Challenge}{Eliezer Yudkowsky}{2008}

Some prophets say that machines will do all the work, and we will never have to lift a finger again. Won't we be bored? No, we will play computer games. But games contain tasks that are ``effortful,'' so we will write programs to help with them. But why have one programmer make a game harder and another make it easier? Make it easy from the start, so you get more gold per hour, until all that is left is a screen that says ``YOU WIN,'' forever, and perhaps a robot to watch that too. That is like the Christian Heaven in one respect, though not supernatural, since it could be built: eternal laziness ``sounds like good news'' to someone who still has to work. \nb{Keynes made the same comparison in 1930. He quoted an old charwoman's epitaph, ``For I'm going to do nothing for ever and ever,'' and wrote ``This was her heaven.''}

And what is a computer game ``except synthetic work''? I ask whether that is a wasted step. Games reduce stress and increase engagement but cost artificiality and isolation. \nb{Bernard Suits, in \textsc{The Grasshopper} (1978), defined playing a game as ``the voluntary attempt to overcome unnecessary obstacles'' and pictured a Utopia without necessary work in which playing games would be what people do.} ``I sometimes think'' we should aim at ``removing low-quality work to make way for high-quality work.'' \nb{In a 1989 study by Csikszentmihalyi and LeFevre, ``the great majority of flow experiences are reported when working, not when in leisure.''}

Some goals cannot be a long-run meaning of life because you can achieve them and be done. It should be made of goals ``good to pursue and not just good to satisfy'': ``valuations that are over 4D states, rather than 3D states,'' games that are fun to play and not only to win. As Timothy Ferriss put it, the question is not ``What do I want?'' but ``What would excite me?'' Some goals, like curing cancer, are worth only their result, since one patient's suffering outweighs the fun of solving it: if after twenty years of your own work an alien offers a cure for thirty bucks, you take it. Such goals are 3D predicates, false now and wanted true later. You want to be there, not go there. They are worth pursuing now, but are not ``long-run fun'' or plausible goals of galactic civilizations.

Why not passive experiencing, then, like the Buddhist Heaven? I am not sure how to build a passive mind. Even lying in bed, your thoughts come from brain areas built to solve problems, and how much brain could you remove and keep the experience of pleasure? I will not touch that. My simpler answer is that ``I wouldn't actually prefer to be a passive experiencer.'' Without Buddha saying Nirvana is the end of existence, it seems to sound like good news only on first hearing. Natural selection built my mind to do things. ``Because it's human nature'' justifies nothing; there is human nature, what we are, and humane nature, what, being human, we wish we were. But ``I don't want to change my nature toward a more passive object,'' and that ``is a justification.'' ``A happy blob is not what, being human, I wish to become.'' \nb{Robert Nozick had used the same image against his experience machine: we want to be ``a certain sort of person,'' not an ``indeterminate blob'' in a tank.}

As I argued earlier about love, what we value can be the real thing and not only the experience of it: it can matter whether your lover is a real person or a realistic nonsentient chatbot, even if you never know. Likewise real challenge and real effort.

In a race, the other racers must be real and I must be able to lose, not in the sense of physical determinism but in that no outside optimizer chose my victory; I must win by my own skill and will, not by pressing a button (though I did not design my legs, so a race of robot cars is a purer contest of their designers, and there is room to improve the human condition). And a sentient being must experience it, not a nonsentient skeleton imitation of the race run trillions of times a second. ``There must be the true effort, the true victory, and the true experience'': the journey, the destination and the traveler.
